\documentclass[11pt]{article}
\usepackage[utf8]{inputenc}
\usepackage{amsmath,amssymb,amsthm,mathtools}
\usepackage{booktabs}
\usepackage[margin=0.9in]{geometry}
\usepackage{natbib}
\usepackage{xcolor}
\usepackage[colorlinks=true,linkcolor=blue!60!black,citecolor=blue!60!black,urlcolor=blue!60!black]{hyperref}
\hypersetup{pdftitle={Minibatch Heavy-Ball SGD with Mode-Dependent Kurtosis: Threshold, Exact Rate, Prefactor and Critical Slowing Down},pdfauthor={Gabriel Paiva}}
\usepackage[expansion=false]{microtype}

\newtheorem{theorem}{Theorem}
\newtheorem{proposition}[theorem]{Proposition}
\newtheorem{lemma}[theorem]{Lemma}
\newtheorem{corollary}[theorem]{Corollary}
\theoremstyle{remark}
\newtheorem{remark}[theorem]{Remark}

\newcommand{\E}{\mathbb{E}}
\newcommand{\R}{\mathbb{R}}
\newcommand{\tr}{\operatorname{tr}}
\newcommand{\T}{\mathcal T}
\newcommand{\V}{\mathcal V}

\title{Minibatch Heavy-Ball SGD with Mode-Dependent Kurtosis:\\ Threshold, Exact Rate, Prefactor and Critical Slowing Down}
\author{Gabriel Paiva\\ \texttt{gabriel@echdesign.co.uk}}
\date{November 2026}

\begin{document}
\maketitle

\begin{abstract}
For heavy-ball SGD on a quadratic in $\R^P$ with minibatch Hessians $H_{B_t}$ of mean $H\succ0$, we study the expected loss at finite dimension $P$ and finite batch $B$ when the features have independent coordinates in the eigenbasis of $H$, with a fourth moment that may differ from mode to mode, possibly multiplied by a common random scale (per-coordinate Bernoulli gates, a per-sample gate on an arbitrary $H$, Student-$t$ features). When the fourth moments differ across modes the curvature noise is not of the spectrally expressible form of earlier generating-function analyses. We show that the second-moment operator $\T$ has the spectral radius of its restriction to diagonal blocks at every step size, that the mean-square threshold is the root of a rank-one secular equation, that twice the expected loss obeys an exact discrete renewal equation, and that $\rho(\T)=1/z_\star$, where $z_\star$ is the smallest positive root of $z\widehat K(z)=1$ for the generating function $\widehat K$ of the renewal kernel. When $\gamma(2+\kappa_k)\ge1$ for every mode $k$ (for instance $\kappa_k\ge-1$), where $\gamma$ is the fourth moment of the random scale and $3+\kappa_k$ that of mode $k$, the loss approaches its noise floor as a single exponential $C\rho(\T)^t$ with an explicit $C>0$, and $1-\rho(\T)$ vanishes linearly at the threshold with an explicit slope. For Gaussian features the secular equation \citep{ziyin2022strength,wang2026momentum}, the renewal equation \citep{wang2026momentum} and the rate as a root of the same equation, with a prefactor in the divergent regime \citep{velikanov2023view}, are known; for isotropic gated features the threshold and the renewal equation are in \citet{everett2026sparse}; the general threshold is in \citet{cai2026llm}. Consequences: a one-direction stability test can overestimate the stable step by two orders of magnitude while the true threshold changes by about $1\%$ between Gaussian and Rademacher features. In the EMA form of momentum the secular sum depends on $\beta$ only through $(1-\beta)/(1+\beta)$, so the threshold and every level set of $S$ increase with $\beta$, by an amount between $\frac{\beta}{1+\beta}\lambda_{\min}$ and $\frac{\beta}{1+\beta}\lambda_{\max}$ in $1/\eta_{\mathrm{EMA}}$, with a finite limit as $\beta\to1$ for $P\ge2$, and the ratio of the threshold to the deterministic one depends on $(\beta,B)$ only through $B(1-\beta)/(1+\beta)$. With label noise a fixed loss floor is a level set of $S$, the rate satisfies $1-\rho(\T)\le1-S$, and the speedup of momentum at a fixed floor ranges from none (one direction, high floor) to above the factor $2$ found for a single isotropic feature (anisotropic $H$, moderate floor). In the general setting, at fixed effective step $\eta/(1-\beta)$ momentum strictly raises the threshold and lowers $\tr(\Sigma_\xi\Sigma_0)$, but can raise the loss.
\end{abstract}

\section{Introduction}

Consider heavy-ball SGD on a quadratic loss in $\R^P$ with mean Hessian $H=\sum_{k=1}^P\lambda_ku_ku_k^\top\succ0$, $\lambda_1\ge\dots\ge\lambda_P>0$:
\begin{equation}\label{eq:hb}
w_{t+1}=\big((1+\beta)I-\eta H_{B_t}\big)w_t-\beta w_{t-1}-\eta\xi_t,\qquad\beta\in[0,1),\ \eta>0,
\end{equation}
where $(H_{B_t},\xi_t)$ are i.i.d.\ across steps and independent of $(w_0,w_{-1})$, with $\E H_{B_t}=H$, $\E\|H_{B_t}\|_F^2<\infty$, $\E\xi_t=0$ and $\Sigma_\xi=\E\xi_t\xi_t^\top\succ0$. Nothing is assumed on the joint law of $H_{B_t}$ and $\xi_t$ within a step. In the EMA form of momentum, $m_{t+1}=\beta m_t+(1-\beta)g_t$, $w_{t+1}=w_t-\eta_{\mathrm{EMA}}m_{t+1}$, the iterates are those of \eqref{eq:hb} with $\eta=(1-\beta)\eta_{\mathrm{EMA}}$. Write $s_t=(w_t,w_{t-1})$, $Q_t=\E s_ts_t^\top$, $F_t=\begin{pmatrix}(1+\beta)I-\eta H_{B_t}&-\beta I\\ I&0\end{pmatrix}$, $\bar F=\E F_t$, $\V(X)=\E[(H_{B_t}-H)X(H_{B_t}-H)]$ and $\T(Q)=\E[F_tQF_t^\top]$.

\paragraph{What is known.}
\emph{Stability.} With non-degenerate additive noise, bounded second moments are equivalent to $\rho(\T)<1$ \citep[Ch.~3]{costa2005mjls}. For LMS with Gaussian inputs ($B=1$, $\beta=0$) the recursion of the diagonal second moments is a diagonal plus rank-one matrix, and its eigenvalues are the roots of the corresponding secular equation \citep[eqs.~(20) and (23)]{feuer1985convergence}; for momentum LMS the characteristic polynomial has a factor $Q(z)$ of the form $1-z\widehat K(z)$ below \citep[Section~IV]{roy1990momentum}. For SGD at interpolating minima the exact threshold is \citet[Theorem~5]{mulayoff2024exact}; for heavy ball and arbitrary non-commuting minibatch Hessians it is \citet[Theorem~F.1]{cai2026llm} (Proposition~\ref{prop:cai} below), whose Lemma~F.13 contains the operator of Theorem~\ref{thm:lyap}(2). The stationary covariance equation with momentum, under ergodicity, is due to \citet{liu2021noise} (Theorem~1 of \citet{ziyin2022strength}); for Gaussian features and label noise \citet[Theorem~2]{ziyin2022strength} solve it in closed form, and its denominator is the secular equation of Theorem~\ref{thm:main} at $\kappa_k\equiv0$.

\emph{Loss dynamics.} Without momentum, \citet[Theorem~1]{bordelon2022learning} give the exact loss for Gaussian features through the matrix $\T_{\mathrm{diag}}$ below at $\beta=0$, and their Theorem~2 gives it for arbitrary features as the iterate of the $P^2$-dimensional second-moment recursion, whose coefficients are the fourth-moment tensor (their eq.~(48)), without additive noise and without a threshold or rate formula. For Gaussian features with momentum and batch, \citet{velikanov2023view} write the loss generating function in closed form (their spectrally expressible approximation is exact there), state stability as $\widetilde U(1)<1$, and in the divergent regime obtain the rate from $r\widetilde U(r)=1$ and a prefactor under an exponential ansatz (their eq.~(24)); \citet[Appendix~B]{wang2026momentum} write the loss as a renewal equation, obtain the threshold from its kernel mass (their eq.~(56)) and bound it (their eq.~(57)). For isotropic least squares with Bernoulli-gated Gaussian inputs, \citet{everett2026sparse} give an exact discrete renewal equation at every $d$ and $B$ (their Lemma~C.1), the threshold (their eq.~(6)), the floor and geometric convergence (their Corollary~C.2) and the per-step spectrum through a cubic (their Lemma~H.4); their vocabulary model has rows that decouple (their Lemma~H.12), they note that an anisotropic $H$ gives one coupled block per eigenvalue (their Remark~B.2), and they propose anisotropic features as future work. In high-dimensional limits the loss of SGD obeys a Volterra equation whose Malthusian root gives the rate, with an explicit prefactor \citep[Theorem~E.2]{paquette2021sgd}; with momentum, \citet{paquette2021momentum} find that small-batch heavy ball gives no improvement over SGD once the step is adjusted, and \citet{lee2022minibatch} study the interpolating case.

\emph{One-direction tests.} The curvature variance along the top eigenvector gives a necessary condition for stability \citep[Proposition~9]{mulayoff2024exact}; batch sharpness above $2/\eta$ certifies instability, and the criterion is one-sided by design \citep[Theorem~1]{andreyev2024eoss}; \citet[eq.~(10)]{andreyev2026momentum} give the one-dimensional threshold with momentum.

\paragraph{What this note adds.}
\begin{enumerate}
\item A model beyond the spectrally expressible class: independent coordinates with mode-dependent fourth moments and a common random scale. For it, $\rho(\T)=\rho(\T_{\mathrm{diag}})$ at every $\eta$, the secular threshold and the exact renewal equation (Lemma~\ref{lem:V}, Theorem~\ref{thm:main}(a),(b)). It contains the Gaussian case and the isotropic gated case of \citet{everett2026sparse}, and extends the latter to an arbitrary $H$ (Corollary~\ref{cor:special}).
\item A proof that $\rho(\T)=1/z_\star$ at finite $P$ for every $\eta$ (Theorem~\ref{thm:main}(c)); on the stable side with additive noise, the single-exponential approach to the floor with an explicit $C>0$, and the slope of critical slowing down at the threshold (Theorem~\ref{thm:main}(d),(e)). We did not find the last two in prior work, also for Gaussian features.
\item The divergence half of the stability criterion, with linear growth from every initialization and without any assumption on the joint law of $(H_{B_t},\xi_t)$ (Lemma~\ref{lem:stab}; the boundedness half with correlated forcing is \citet[Proposition~C.3]{cai2026llm}), and the step-covariance identity $\E[\Delta w\Delta w^\top]=\frac{\eta}{1+\beta}(H\Sigma_0+\Sigma_0H)$ (Theorem~\ref{thm:lyap}(1)).
\item Momentum and batch in EMA form (Corollary~\ref{cor:special}(iv)): monotonicity of every level set of $S$ in $\beta$ with two-sided bounds, the limit $\beta\to1$, and the dependence of $\eta^\star/\eta_{\mathrm{det}}$ on $B(1-\beta)/(1+\beta)$ alone; with label noise, the rate at a fixed loss floor (Proposition~\ref{prop:floor}), extending the isotropic comparisons of \citet[Proposition~H.10, Lemma~C.5]{everett2026sparse}.
\item The size of the gap between the one-direction threshold and the true threshold for non-Gaussian features, against a near-universality of the true threshold in the kurtosis (Proposition~\ref{prop:1d}; for Gaussian features \citet[eq.~(57)]{wang2026momentum} give the same upper bound and the lower bounds $\alpha\lambda_1$ and $p\kappa_\beta\tr H$), and, in the general setting, the strict increase of $\eta^\star/(1-\beta)$ in $\beta$ with an explicit margin and a signed finite-step comparison at fixed effective step (Proposition~\ref{prop:eff}; the monotonicity is eq.~(67) of \citet{cai2026llm}, written there as a non-strict inequality, and the margin is new).
\end{enumerate}
Proofs are in Appendix~\ref{app:proofs}; every number is printed by \texttt{paper/checks/note-v2-checks.py}.

\section{General facts}\label{sec:general}

\begin{lemma}[Stability without assumptions on the joint law]\label{lem:stab}
Assume the setting of \eqref{eq:hb} and that $(w_0,w_{-1})$ has finite second moments. Let $G=\operatorname{blkdiag}(\eta^2\Sigma_\xi,0)$.
\begin{enumerate}
\item[(i)] $\T$ preserves the symmetric and the positive semidefinite matrices, $\|\T^k\|=\|\T^k(I)\|$ in the operator norm induced by the spectral norm, and $\rho(\T)=\rho(\T|_{\mathbb S^{2P}})$.
\item[(ii)] If $\rho(\T)<1$, then $Q_t\to Q_\infty=\sum_{k\ge0}\T^k(G)$ and $\E w_t\to0$, whatever the joint law of $(H_{B_t},\xi_t)$.
\item[(iii)] If $\rho(\T)\ge1$, there is $c>0$ with $\E\|w_t\|^2+\E\|w_{t-1}\|^2\ge ct$ for every $t\ge1$.
\end{enumerate}
The assumption $H\succ0$ cannot be dropped in (iii): for $P=1$, $H=0$, $H_{B_t}=\epsilon_t\in\{\pm1\}$ equiprobable, $\xi_t=-\epsilon_t$ and $w_0=w_{-1}=1$, one has $w_t\equiv1$ while $\rho(\T)\ge\rho(\bar F)^2=1$.
\end{lemma}
Part (ii), with a forcing correlated with the transition, is \citet[Proposition~C.3]{cai2026llm}, and their Proposition~A.3 gives, above their threshold, an unbounded second moment from some deterministic initialization; (iii) gives linear growth from every initialization.

\begin{theorem}[Stationary covariance equation]\label{thm:lyap}
Let $\rho(\T)<1$ and let $\Sigma_0=\E w_tw_t^\top$, $D=\E[(w_{t+1}-w_t)(w_{t+1}-w_t)^\top]$ in the stationary regime of Lemma~\ref{lem:stab}(ii). Then
\begin{enumerate}
\item[(1)] $D=\frac{\eta}{1+\beta}(H\Sigma_0+\Sigma_0H)$;
\item[(2)] $\displaystyle(H\Sigma_0+\Sigma_0H)-\frac{\eta}{1+\beta}H\Sigma_0H+\frac{\eta\beta}{(1-\beta)(1-\beta^2)}[H,[H,\Sigma_0]]-\frac{\eta}{1-\beta}\V(\Sigma_0)=\frac{\eta}{1-\beta}\Sigma_\xi.$
\end{enumerate}
\end{theorem}
With $\V$ replaced by a given noise covariance, (2) is the equation of \citet{liu2021noise} (Theorem~1 of \citet{ziyin2022strength}), written with $H^2\Sigma_0+\Sigma_0H^2=[H,[H,\Sigma_0]]+2H\Sigma_0H$. Write (2) as $(\mathcal L_0-\eta\mathcal B_\beta)\Sigma_0=\frac{\eta}{1-\beta}\Sigma_\xi$ with $\mathcal L_0X=HX+XH$ and $\mathcal B_\beta X=\frac{1}{1+\beta}HXH-\frac{\beta}{(1-\beta)(1-\beta^2)}[H,[H,X]]+\frac1{1-\beta}\V(X)$, both self-adjoint on $(\mathbb S^P,\langle X,Y\rangle=\tr XY)$.

\begin{proposition}[\citealp{cai2026llm}, Theorem~F.1]\label{prop:cai}
$\rho(\T)<1$ if and only if $\mathcal L_0-\eta\mathcal B_\beta\succ0$ on $\mathbb S^P$, i.e.\ $\eta<\eta^\star\triangleq1/\lambda_{\max}(\mathcal L_0^{-1/2}\mathcal B_\beta\mathcal L_0^{-1/2})$.
\end{proposition}
\citet{cai2026llm} state it as $(1-b)/\alpha>r_C(b)$ for a quotient $q_{C,b}$ maximized over positive semidefinite $X$. The maximum over all of $\mathbb S^P$ is the same: in their notation $D^{-1}\V$ preserves the positive semidefinite cone \citep[Lemma~F.13]{cai2026llm}, so by the Krein--Rutman theorem its spectral radius is an eigenvalue with a positive semidefinite eigenmatrix, and $D^{-1}\V$ is similar to the self-adjoint $D^{-1/2}\V D^{-1/2}$, whose top eigenvalue is therefore that spectral radius. At $\beta=0$, under the sampling and positive-semidefinite sample-Hessian assumptions of \citet[Theorem~5]{mulayoff2024exact}, it recovers their threshold.

\section{Independent coordinates with mode-dependent fourth moments}\label{sec:rf}

Let $H_{B_t}=\frac1B\sum_{b=1}^Bz_{t,b}z_{t,b}^\top$ with i.i.d.\ $z_{t,b}=\zeta_{t,b}H^{1/2}\omega_{t,b}$, where the scalar $\zeta$ is independent of $\omega$, $\E\zeta^2=1$, $\gamma\triangleq\E\zeta^4$ ($\ge1$), and the coordinates $\omega_k=u_k^\top\omega$ are independent with $\E\omega_k=0$, $\E\omega_k^2=1$, $\E\omega_k^4=3+\kappa_k$. Put $p=1/B$, $\kappa_\beta=\frac{1+\beta}{1-\beta}$, and for each mode
\[
c_k\triangleq p\big(\gamma(2+\kappa_k)-1\big)\lambda_k^2,\qquad \alpha_k\triangleq1+\kappa_\beta\frac{c_k}{\lambda_k^2},\qquad x_k\triangleq\frac{\eta\lambda_k}{2(1+\beta)},\qquad a_k\triangleq1+\beta-\eta\lambda_k .
\]
Examples: (a) $\zeta\equiv1$, $\kappa_k\equiv\kappa$ (Gaussian $\kappa=0$, Rademacher $\kappa=-2$); (b) per-coordinate gates $\omega_k=(s_k/\sqrt{q_k})\tilde\omega_k$, $s_k\sim\mathrm{Bernoulli}(q_k)$ independent, so $\kappa_k=(3+\tilde\kappa_k)/q_k-3$ and $\lambda_k$ carries the frequency $q_k$; (c) the per-sample gate $\zeta=s/\sqrt q$, $s\sim\mathrm{Bernoulli}(q)$, of \citet{everett2026sparse}, now on an arbitrary $H$: $\gamma=1/q$; (d) Student-$t_\nu$ features, $\nu>4$: $\gamma=\frac{\nu-2}{\nu-4}$, $\kappa_k=0$.

\begin{lemma}[Curvature noise and diagonal reduction]\label{lem:V}
In the eigenbasis of $H$, for $X\in\mathbb S^P$,
\[
\V(X)_{kk}=p\gamma\lambda_k\tr(HX)+c_kX_{kk},\qquad \V(X)_{jk}=p(2\gamma-1)\lambda_j\lambda_kX_{jk}\quad(j\ne k).
\]
Hence $\T$ maps the space $\mathbb S^{2P}_{\mathrm{diag}}$ of matrices whose four $P\times P$ blocks are diagonal in this basis into itself, and $\rho(\T)=\rho(\T_{\mathrm{diag}})$ for every $\eta>0$, where $\T_{\mathrm{diag}}=\T|_{\mathbb S^{2P}_{\mathrm{diag}}}$.
\end{lemma}

In coordinates $(\E w_k^2,\E w_kw'_k,\E w_k'^2)_{k\le P}$ ($w'=w_{t-1}$), $\T_{\mathrm{diag}}=\operatorname{blkdiag}(\widetilde B_k)+\eta^2p\gamma\,rr^\top$ with $r=(\lambda_ke_1)_{k\le P}$ and
\[
\widetilde B_k=\begin{pmatrix}a_k^2+\eta^2c_k&-2a_k\beta&\beta^2\\ a_k&-\beta&0\\ 1&0&0\end{pmatrix}.
\]
Let $g_k$ be the impulse response of mode $k$ of deterministic heavy ball ($g_k(0)=1$, $g_k(1)=a_k$, $g_k(\tau+1)=a_kg_k(\tau)-\beta g_k(\tau-1)$), $\psi_k(\tau)=\eta^2g_k(\tau)^2$, $h_k(\tau)=\psi_k(\tau)+c_k\sum_{s<\tau}\psi_k(\tau-1-s)h_k(s)$, and
\[
\begin{gathered}
K_\tau=p\gamma\sum_k\lambda_k^2h_k(\tau),\qquad \widehat\psi_k(z)=\frac{\eta^2(1+\beta z)}{(1-\beta z)\big(1-(a_k^2-2\beta)z+\beta^2z^2\big)},\\ \widehat h_k=\frac{\widehat\psi_k}{1-c_kz\widehat\psi_k},\qquad \widehat K=p\gamma\sum_k\lambda_k^2\widehat h_k .
\end{gathered}
\]
$\widehat\psi_k$, $\widehat h_k$, $\widehat K$ are the generating functions of $\psi_k$, $h_k$, $K$ (rational in $z$). Let $\sigma_k^2=u_k^\top\Sigma_\xi u_k$ and
\[
S(\eta)\triangleq\frac{\eta p\gamma}{2(1-\beta)}\sum_{k=1}^P\frac{\lambda_k}{1-\alpha_kx_k}.
\]

\begin{theorem}[Threshold, renewal equation, rate, prefactor, slowing down]\label{thm:main}
In this model:
\begin{enumerate}
\item[(a)] \textbf{Threshold.} $\rho(\T)<1$ if and only if $\eta<\eta^\star$, where $\eta^\star$ is the unique root of $S(\eta)=1$ with $\eta\max_k\max(\alpha_k,0)\lambda_k<2(1+\beta)$; for $\eta<\eta^\star$, $\ell_\infty\triangleq\tr(H\Sigma_0)=f_\infty/(1-S(\eta))$ with $f_\infty=\sum_k\frac{\eta\sigma_k^2}{2(1-\beta)(1-\alpha_kx_k)}$.
\item[(b)] \textbf{Renewal equation.} From $w_0=w_{-1}=0$, $\ell_t=\tr(H\E w_tw_t^\top)$ satisfies $\ell_t=f_t+\sum_{s<t}K_{t-1-s}\ell_s$ with $f_t=\sum_{\tau<t}\phi_\tau$, $\phi_\tau=\sum_k\lambda_k\sigma_k^2h_k(\tau)$; equivalently $\widehat\ell(z)=z\widehat\phi(z)/\big((1-z)(1-z\widehat K(z))\big)$.
\item[(c)] \textbf{Rate.} If $\rho(\T)>0$ (automatic for $\beta>0$, since $\rho(\T)\ge\rho(\bar F)^2\ge\beta$), then $\rho(\T)=1/z_\star$, where $z_\star$ is the smallest positive root of $z\widehat K(z)=1$, $\widehat K$ being read as the rational function above (for $c_k<0$ its power series may diverge at $z_\star$).
\item[(d)] \textbf{Single exponential.} If $c_k\ge0$ for every $k$ (i.e.\ $\gamma(2+\kappa_k)\ge1$) and $\eta<\eta^\star$, then $z_\star>1$ is the only root of $z\widehat K(z)=1$ in $|z|\le z_\star$, it is simple, and
\[
\ell_\infty-\ell_t=C\rho(\T)^t+O(r^{-t})\ \text{ for some } r>z_\star,\qquad C=\frac{\widehat\phi(z_\star)}{(z_\star-1)\,\frac{d}{dz}\big(z\widehat K(z)\big)\big|_{z_\star}}>0 .
\]
\item[(e)] \textbf{Critical slowing down.} Under the condition of (d), as $\eta\uparrow\eta^\star$,
\[
1-\rho(\T)=\frac{S'(\eta^\star)}{\sum_{\tau\ge0}(\tau+1)K_\tau}\,(\eta^\star-\eta)+O\big((\eta^\star-\eta)^2\big),\qquad S'(\eta^\star)>0,
\]
the kernel being taken at $\eta^\star$, where $\sum_\tau K_\tau=S(\eta^\star)=1$; the denominator is the mean renewal time.
\end{enumerate}
\end{theorem}

The condition $c_k\ge0$ is $\kappa_k\ge1/\gamma-2$; it holds for every $\kappa_k\ge-1$ (as $\gamma\ge1$), and for Rademacher modes ($\kappa_k=-2$) it fails for every $\gamma$. For $c_k<0$ the kernel can have both signs; (a)--(c) still hold, while $h_k$ need not be summable (for $P=1$, $\lambda_1=1$, $\beta=0$, $B=1$, Rademacher features and $\eta=1.5$, $\widehat h_1=2.25/(1+2z)$, $h_1$ grows geometrically and $\rho(\T)=1/4=1/z_\star$), and the single exponential of (d) can fail: for $P=1$ Rademacher features $H_{B_t}\equiv\lambda_1$, $\T_{\mathrm{diag}}=B_1$, all eigenvalues of $B_1$ have modulus $\beta$ in the underdamped regime, and $(\ell_\infty-\ell_t)/\beta^t$ does not converge. From a general initial condition with $\E w_0=\E w_{-1}=0$ (or with $\E[\xi_t(H_{B_t}-H)_{jk}]=0$ for all $j,k$), (b) holds with $f_t$ replaced by $f_t+\sum_k\lambda_k\,\epsilon_k(t)$, where $\epsilon_k(t)$ propagates the initial $2\times2$ block of mode $k$ through $\widetilde B_k$; otherwise a cross term $\eta^2\E[(H_{B_t}-H)\bar w_t\xi_t^\top]$ enters; without additive noise this is the interpolating setting of \citet{lee2022minibatch} and \citet{everett2026sparse}.

\begin{corollary}[Special cases]\label{cor:special}
\begin{enumerate}
\item[(i)] Gaussian features ($\gamma=1$, $\kappa_k\equiv0$): $S(\eta)=1$ is the vanishing of the denominator $1-\frac1B\tr[\Lambda AG_\mu^{-1}]$ of \citet[Theorem~2]{ziyin2022strength} (scalar learning rate) and the critical equation (56) of \citet{wang2026momentum}.
\item[(ii)] Per-sample gate on an arbitrary $H=q\tilde H$ ($\gamma=1/q$, $\kappa_k=0$, $\tilde H=\operatorname{diag}(\tilde\lambda_k)$), in EMA form:
\[
\frac{\eta_{\mathrm{EMA}}}{2B}\sum_{k=1}^P\frac{\tilde\lambda_k}{1-\frac{\eta_{\mathrm{EMA}}\tilde\lambda_k}{2}\Big(q\frac{1-\beta}{1+\beta}+\frac{2-q}{B}\Big)}=1 .
\]
For $\tilde H=I_d$ this is $1/\eta_{\mathrm{EMA}}^\star=\frac{d+2-q}{2B}+\frac q2\frac{1-\beta}{1+\beta}$, equation~(6) of \citet{everett2026sparse}, and the rate equation of (c) reduces to a cubic; their Lemma~H.4 gives the same per-step spectrum, written as $1+\varepsilon z$ at $\beta=1-\varepsilon$.
\item[(iii)] Per-coordinate gates: as $q_k\to0$ at fixed $\tilde\lambda_k$, $\alpha_k\lambda_k\to p\kappa_\beta(3+\tilde\kappa_k)\tilde\lambda_k$, so in the limit of rare coordinates the admissible range of (a) becomes $\eta_{\mathrm{EMA}}<2B/((3+\tilde\kappa_k)\tilde\lambda_k)$, which in EMA form does not depend on $\beta$ (at finite $q_k$ the range is $\eta_{\mathrm{EMA}}\tilde\lambda_k\big(q_k\frac{1-\beta}{1+\beta}+p(3+\tilde\kappa_k-2q_k)\big)<2$, which does); all coordinates are coupled through the trace term of $S$. The gates here are independent across coordinates, unlike the exclusive tokens of \citet{everett2026sparse}, whose rows decouple and whose noise ceiling $\eta<2B/(d+2-p)$ \citep[Section~4.2]{everett2026sparse} comes from the trace term, while this one comes from the diagonal self-term.
\item[(iv)] EMA form, momentum and batch. Write $\lambda_{\max}=\lambda_1$, $\lambda_{\min}=\lambda_P$, put $\theta=\frac{1-\beta}{1+\beta}$, $d_k=\frac12\big(\theta\lambda_k+c_k/\lambda_k\big)$ and $M_\infty=\operatorname{diag}\big((\gamma(2+\kappa_k)-1)\lambda_k\big)+\gamma\sqrt\lambda\sqrt\lambda^\top$. With $\eta=(1-\beta)\eta_{\mathrm{EMA}}$,
\[
S=\frac{\eta_{\mathrm{EMA}}\,p\gamma}{2}\sum_{k=1}^P\frac{\lambda_k}{1-\eta_{\mathrm{EMA}}d_k},\qquad \frac1{\eta^\star_{\mathrm{EMA}}}=\lambda_{\max}\Big(\operatorname{diag}(d_k)+\frac{p\gamma}{2}\sqrt\lambda\sqrt\lambda^\top\Big),
\]
and $\beta$ enters only through $\theta$. (a) For every level $s\in(0,1]$ the root $\eta_{\mathrm{EMA}}(\beta;s)$ of $S=s$ in the admissible range ($s=1$ gives $\eta^\star_{\mathrm{EMA}}$) increases strictly with $\beta$, and
\[
\frac{\beta}{1+\beta}\lambda_{\min}\le\frac1{\eta_{\mathrm{EMA}}(0;s)}-\frac1{\eta_{\mathrm{EMA}}(\beta;s)}\le\frac{\beta}{1+\beta}\lambda_{\max}.
\]
(b) As $\beta\uparrow1$, $1/\eta^\star_{\mathrm{EMA}}$ decreases to $\frac p2\lambda_{\max}(M_\infty)$, and $\lambda_{\max}(M_\infty)\ge\sqrt{\lambda_{(1)}\lambda_{(2)}}>0$ for $P\ge2$ ($\lambda_{(1)}\ge\lambda_{(2)}$ the two largest eigenvalues), so the EMA threshold stays bounded; for $s<1$, $\eta_{\mathrm{EMA}}(\beta;s)$ increases to the finite root of $\frac{xp\gamma}{2}\sum_k\lambda_k/(1-xc_k/(2\lambda_k))=s$. (c) With $\eta_{\mathrm{det}}=2(1+\beta)/\lambda_{\max}$, $\eta^\star/\eta_{\mathrm{det}}=\lambda_{\max}/\lambda_{\max}\big(H+\frac{1+\beta}{(1-\beta)B}M_\infty\big)$ depends on $(\beta,B)$ only through $B(1-\beta)/(1+\beta)$. If $\lambda_{\max}(M_\infty)>0$, then $\eta^\star\ge\eta_{\mathrm{det}}/2$ if and only if $B\ge B_{1/2}(\beta)=\frac{1+\beta}{1-\beta}\lambda_{\max}(G^{-1/2}M_\infty G^{-1/2})$, $G=2\lambda_{\max}I-H$ (in the eigenbasis of $H$), and $\lambda_{\max}(M_\infty)/(2\lambda_{\max})\le B_{1/2}(0)\le\lambda_{\max}(M_\infty)/\lambda_{\max}$.
For $\tilde H=I_d$ in (ii), (a)--(c) can be read from eq.~(6) of \citet{everett2026sparse}; for spectrally expressible curvature noise, Gaussian features among them, the limit $\beta\to1$ of the effective threshold is in \citet[eq.~(23)]{velikanov2023view}; mode-dependent $\kappa_k$ is not of that form. Their critical batch size \citep[Proposition~2]{everett2026sparse} counts samples and steps to a target loss and is a different quantity from $B_{1/2}$.
\end{enumerate}
\end{corollary}

\section{Two consequences}\label{sec:cons}

\begin{proposition}[One-direction test against the true threshold]\label{prop:1d}
Let $\gamma=1$, $\kappa_k\equiv\kappa$, $\alpha=1+p\kappa_\beta(1+\kappa)$, $\mu^\star=2(1+\beta)/\eta^\star$ and $\mu_{\mathrm{1D}}=\lambda_1\big(1+p\kappa_\beta(2+\kappa)\big)$, so that $2(1+\beta)/\mu_{\mathrm{1D}}$ is the one-dimensional threshold along $u_1$ \citep[eq.~(10)]{andreyev2026momentum} with $\mathrm{Var}(u_1^\top H_{B_t}u_1)=p(2+\kappa)\lambda_1^2$. Then $\mu^\star\ge\mu_{\mathrm{1D}}$ and
\[
\alpha\,\frac{\tr H^2}{\tr H}\le\mu^\star-p\kappa_\beta\tr H\le\max(\alpha,0)\lambda_1 .
\]
In particular, at high effective rank $r_{\mathrm{eff}}=\tr H/\lambda_1$ the true threshold depends on $\kappa$ only through a relative term $|(p\kappa_\beta)^{-1}+1+\kappa|/r_{\mathrm{eff}}$, while for Rademacher features ($\mu_{\mathrm{1D}}=\lambda_1$; $\alpha<0$ when $B<\kappa_\beta$) the ratio of the one-dimensional threshold to the true one is at least $p\kappa_\beta r_{\mathrm{eff}}+\alpha\tr H^2/(\lambda_1\tr H)$.
\end{proposition}

\begin{table}[!ht]
\centering
\caption{$P=256$, $\lambda_k=k^{-1/2}$ ($r_{\mathrm{eff}}=30.57$), $\beta=0.9$, $B=4$.}
\label{tab:1d}
\smallskip
{\small\begin{tabular}{lrrrrr}
\toprule
$\kappa$ & $-2$ (Rademacher) & $-1$ & $0$ (Gaussian) & $2$ & $10$\\
\midrule
true threshold $\eta^\star$ & 0.026304 & 0.026133 & 0.025961 & 0.025617 & 0.024223\\
one-dimensional threshold & 3.8000 & 0.66087 & 0.36190 & 0.19000 & 0.065517\\
ratio & 144.46 & 25.29 & 13.94 & 7.42 & 2.70\\
\bottomrule
\end{tabular}}
\end{table}

Table~\ref{tab:1d}: from Gaussian to Rademacher features the true threshold moves by $1.32\%$ and the one-dimensional threshold by a factor $10.5$; the bound of Proposition~\ref{prop:1d} gives a ratio of at least $144.460$ at $\kappa=-2$ (actual $144.464$). When $1+\kappa$ is of order $r_{\mathrm{eff}}$ or larger, the term $\alpha\lambda_1$ dominates and the kurtosis enters the true threshold at first order.

\begin{proposition}[Fixed effective step]\label{prop:eff}
In the general setting of Section~\ref{sec:general}, let $\V$ be independent of $\beta$. (i) The effective threshold $\eta^\star(\beta)/(1-\beta)$ increases strictly with $\beta\in[0,1)$, and $\frac1{\eta^\star(0)}-\frac{1-\beta}{\eta^\star(\beta)}\ge\frac{\beta}{1+\beta}\lambda_{\min}(H)$. (ii) Fix $\eta_{\mathrm{eff}}>0$ and let $\eta=(1-\beta)\eta_{\mathrm{eff}}$. If $0\le\beta_1<\beta_2<1$ are both stable, then $\tr(\Sigma_\xi\Sigma_0(\beta_2))<\tr(\Sigma_\xi\Sigma_0(\beta_1))$. If $\Sigma_\xi$ is proportional to $H$ the loss $\tr(H\Sigma_0)$ therefore decreases strictly; in general it can increase: for $H=\operatorname{diag}(1,0.3)$, two equiprobable minibatch Hessians $H\pm N$ with $N=\begin{psmallmatrix}-0.6&-0.7\\-0.7&-0.2\end{psmallmatrix}$, $\Sigma_\xi=vv^\top+10^{-3}I$, $v=(-0.8,0.6)$ and $\eta_{\mathrm{eff}}=0.52$, raising $\beta$ from $0.5$ to $0.95$ raises the loss from $0.34072$ to $0.35659$ and lowers $\tr(\Sigma_\xi\Sigma_0)$ from $0.43506$ to $0.34534$.
\end{proposition}

At fixed small step, momentum SGD tracks SGD at step $\eta/(1-\beta)$ \citep{yuan2016momentum}; Proposition~\ref{prop:eff} gives the sign of the departure at finite step, for the noise-weighted second moment.

\begin{proposition}[Rate at a fixed loss floor]\label{prop:floor}
In the model of Section~\ref{sec:rf} let $\Sigma_\xi=\sigma^2H/B$ (least squares with label noise of variance $\sigma^2>0$) and $\eta<\eta^\star$. Then $\ell_\infty=\frac{\sigma^2}{\gamma}\frac{S(\eta)}{1-S(\eta)}$, so a target floor fixes the level $s=\gamma\ell_\infty/(\sigma^2+\gamma\ell_\infty)$ of $S$, and the EMA step that reaches it is $\eta_{\mathrm{EMA}}(\beta;s)$ of Corollary~\ref{cor:special}(iv). If $c_k\ge0$ for every $k$, then, with the mean renewal time $\bar\tau=\frac1S\sum_{\tau\ge0}(\tau+1)K_\tau$,
\[
\rho(\T)\ge S^{1/\bar\tau}\ge S,\qquad\text{so}\qquad 1-\rho(\T)\le1-S=\frac{\sigma^2}{\sigma^2+\gamma\ell_\infty}.
\]
\end{proposition}
The continuum, isotropic analogue of the last bound is $r_c<1-H$ in \citet[Lemma~C.5(c)]{everett2026sparse}. Table~\ref{tab:floor} maximizes $1-\rho(\T)$ over $\beta$ along a fixed floor, with $\rho(\T)=1/z_\star$ computed exactly. In one direction at a high floor momentum does not help; at a low floor the gain approaches $2$ (rows 3 and 5), the factor that \citet[Proposition~H.10]{everett2026sparse} find for isotropic features; at a moderate floor with an anisotropic $H$ it exceeds $2$ (row 4), consistent with the conditioning gain that \citet[Proposition~2]{everett2026sparse} find for a vocabulary at large batch. Numerically, along a fixed floor $(1-\rho(\T))/(1-\beta)\to1-\eta_{\mathrm{EMA}}(\beta;s)/\eta^\star_{\mathrm{EMA}}(\beta)$ as $\beta\uparrow1$ (in the fourth row, $0.61520$ against $0.61516$ at $\beta=1-10^{-4}$); we do not prove this.

\begin{table}[!ht]
\centering
\caption{Rate at a fixed level $s$ of $S$ (fixed floor), maximized over $\beta$ (grid of $318$ values in $[0,1-10^{-4.5}]$, refined by golden-section search). Rows 4--6: $\lambda=(1,0.5,0.2,0.05,0.01)$, $B=8$; row 6 has $\kappa=(5,0,2,8,1)$ and $\gamma=2$.}
\label{tab:floor}
\smallskip
{\small\begin{tabular}{lrrrrr}
\toprule
features & $s$ & $1-\rho$ at $\beta=0$ & best $\beta$ & $1-\rho$ there & gain\\
\midrule
$P=1$, Gaussian, $B=4$ & $0.3$ & 0.53760 & 0 & 0.53760 & 1.000\\
$P=1$, Gaussian, $B=4$ & $0.01$ & 0.14367 & 0.752 & 0.21215 & 1.477\\
$P=1$, Gaussian, $B=4$ & $0.001$ & 0.015825 & 0.970 & 0.028811 & 1.821\\
$P=5$, Gaussian & $0.3$ & 0.023537 & 0.897 & 0.060514 & 2.571\\
$P=5$, Gaussian & $0.01$ & 0.0017510 & 0.996 & 0.0035060 & 2.003\\
$P=5$, mixed $\kappa_k$, $\gamma=2$ & $0.3$ & 0.010833 & 0.942 & 0.016475 & 1.521\\
\bottomrule
\end{tabular}}
\end{table}

\appendix
\section{Proofs}\label{app:proofs}

\paragraph{Lemma~\ref{lem:stab}.}
(i) $\T(Q)^\top=\T(Q^\top)$; $\T^k$ is in Kraus form, hence positive, and the Russo--Dye theorem gives $\|\T^k\|=\|\T^k(I)\|$; Gelfand's formula then gives $\rho(\T)=\lim\|\T^k(I)\|^{1/k}\le\rho(\T|_{\mathbb S^{2P}})$, and the reverse inequality is trivial.

Let $m_t=\E s_t$ and $g_t=(-\eta\xi_t,0)$. Since $s_t$ is independent of $(F_t,g_t)$ and $\E g_t=0$, $m_{t+1}=\bar Fm_t$ and $Q_{t+1}=\T(Q_t)+G+C_t$, with $C_t=\E[(F_t-\bar F)m_tg_t^\top]+(\cdot)^\top$, $\|C_t\|\le c_0\|m_t\|$. Write $\T=\T_0+\Delta$, $\T_0(Q)=\bar FQ\bar F^\top$, $\Delta(Q)=\operatorname{blkdiag}(\eta^2\V(Q_{11}),0)$; both are positive, so $\rho(\bar F)^2=\rho(\T_0)\le\rho(\T)$ \citep{berman1994}.

(ii) If $\rho(\T)<1$, then $\|m_t\|\le c_1r^t$ with $r<1$, and $Q_t=\T^t(Q_0)+\sum_{s<t}\T^{t-1-s}(G+C_s)\to\sum_k\T^k(G)$.

(iii) Let $\rho=\rho(\T)\ge1$. By the Krein--Rutman theorem \citep[Thm.~1.3.2]{berman1994} the adjoint $\T^*(Y)=\E[F_t^\top YF_t]$ has an eigenmatrix $W\succeq0$, $W\ne0$, with eigenvalue $\rho$, and $W_{11}\ne0$ (otherwise $W\succeq0$ forces $W_{12}=0$, and then $F_t^\top WF_t=\operatorname{blkdiag}(W_{22},0)$ forces $W=0$), so $g_0\triangleq\tr(W_{11}\Sigma_\xi)>0$. The centered second moment $\widetilde Q_t=Q_t-m_tm_t^\top$ obeys $\widetilde Q_{t+1}=\T(\widetilde Q_t)+U_t$ with $U_t=\E[\upsilon_t\upsilon_t^\top]\succeq0$, $\upsilon_t=(F_t-\bar F)m_t+g_t$, because the centered state is independent of $(F_t,g_t)$ and has mean zero. With $\bar w_t=\E w_t$ and $\Delta_t=H_{B_t}-H$, $\tr(WU_t)=\eta^2\varphi(\bar w_t)$, $\varphi(a)=\E[(\Delta_ta+\xi_t)^\top W_{11}(\Delta_ta+\xi_t)]$, and
\[
\tr(W\widetilde Q_{t+1})=\rho\tr(W\widetilde Q_t)+\eta^2\varphi(\bar w_t)\ge\tr(W\widetilde Q_t)+\eta^2\varphi(\bar w_t).
\]
The parallelogram identity gives $\varphi(a)+\varphi(-a)=2\E[a^\top\Delta_tW_{11}\Delta_ta]+2g_0\ge2g_0$ for every $a$. Each $2\times2$ block of $\bar F$ in the eigenbasis of $H$ has eigenvalues with product $\beta<1$, and $\det(I-\bar F_k)=\eta\lambda_k>0$; so the eigenvalues of modulus $\ge1$ are semisimple and the only one on the unit circle is $-1$ (a non-real pair on the circle would have product $1$, and $+1$ is excluded by $\det(I-\bar F_k)>0$). If $m_0$ has a component on the eigenvalues of modulus $>1$, $\|m_t\|$ and hence $\tr Q_t$ grow geometrically. Otherwise $\bar w_t=(-1)^ta+o(1)$, so $\varphi(\bar w_{2k})+\varphi(\bar w_{2k+1})\ge g_0$ for large $k$, and $\tr(W\widetilde Q_t)$ grows at least linearly; $\tr Q_t\ge\tr(W\widetilde Q_t)/\|W\|$. Finally $\E\|w_t\|^2+\E\|w_{t-1}\|^2>0$ for every $t\ge1$: if $w_t=w_{t-1}=0$ a.s., then $\eta\xi_{t-1}=-\beta w_{t-2}$ a.s.\ with $w_{t-2}$ independent of $\xi_{t-1}$, so $\xi_{t-1}$ is a.s.\ constant, contradicting $\Sigma_\xi\succ0$. The bound $ct$ follows. The counterexample with $H=0$ satisfies $(H_{B_t}-H)w_t+\xi_t=0$ on the path $w\equiv1$.

\paragraph{Theorem~\ref{thm:lyap}.}
By Lemma~\ref{lem:stab}(ii), $\E w_t=0$ in the stationary regime, and $(H_{B_t}-H,\xi_t)$ is independent of $(w_t,w_{t-1})$, so $\E[(H_{B_t}-H)w_t\xi_t^\top]=0$ whatever the joint law of $(H_{B_t},\xi_t)$. With $A=(1+\beta)I-\eta H$ and $\Sigma_1=\E w_{t+1}w_t^\top$: conditioning on the past gives $\Sigma_1=A\Sigma_0-\beta\Sigma_1^\top$, hence $\Sigma_1=(A\Sigma_0-\beta\Sigma_0A)/(1-\beta^2)$ and $\Sigma_1+\Sigma_1^\top=2\Sigma_0-\frac{\eta}{1+\beta}(H\Sigma_0+\Sigma_0H)$, which is (1) since $D=2\Sigma_0-\Sigma_1-\Sigma_1^\top$. Expanding $\Sigma_0=\E w_{t+1}w_{t+1}^\top$ gives $\Sigma_0=A\Sigma_0A+\beta^2\Sigma_0-\beta(A\Sigma_1+\Sigma_1^\top A)+\eta^2\V(\Sigma_0)+\eta^2\Sigma_\xi$; with $A^2\Sigma_0+\Sigma_0A^2-2\beta A\Sigma_0A=[A,[A,\Sigma_0]]+2(1-\beta)A\Sigma_0A$ and $[A,[A,\Sigma_0]]=\eta^2[H,[H,\Sigma_0]]$ this is $(1-\beta^2)\Sigma_0-\frac{1-\beta}{1+\beta}A\Sigma_0A=-\frac{\beta\eta^2}{1-\beta^2}[H,[H,\Sigma_0]]+\eta^2\V(\Sigma_0)+\eta^2\Sigma_\xi$; expanding $A\Sigma_0A$ and dividing by $\eta(1-\beta)$ gives (2).

\paragraph{Lemma~\ref{lem:V}.}
$\V(X)=p(\E[zz^\top Xzz^\top]-HXH)$ and $\E[z_jz_lz_mz_k]=\gamma\sqrt{\lambda_j\lambda_l\lambda_m\lambda_k}\E[\omega_j\omega_l\omega_m\omega_k]$, which vanishes unless every index appears an even number of times. For $j=k$: $l=m\ne k$ gives $\gamma\lambda_k\lambda_lX_{ll}$ and $l=m=k$ gives $\gamma(3+\kappa_k)\lambda_k^2X_{kk}$. For $j\ne k$: $(l,m)=(j,k)$ and $(k,j)$ each give $\gamma\lambda_j\lambda_kX_{jk}$ on symmetric $X$. Subtracting $(HXH)_{jk}=\lambda_j\lambda_kX_{jk}$ gives the formulas. Thus $\V$ maps diagonal matrices to diagonal ones and off-diagonal entries to themselves, and $\bar F$ acts mode by mode, so $\T(\mathbb S^{2P}_{\mathrm{diag}})\subseteq\mathbb S^{2P}_{\mathrm{diag}}$. Since $I\in\mathbb S^{2P}_{\mathrm{diag}}$, $\|\T^k\|=\|\T^k(I)\|=\|\T_{\mathrm{diag}}^k(I)\|\le\|\T^k_{\mathrm{diag}}\|\le\|\T^k\|$ by Lemma~\ref{lem:stab}(i), and Gelfand's formula gives $\rho(\T)=\rho(\T_{\mathrm{diag}})$.

\paragraph{Theorem~\ref{thm:main}.}
\emph{Generating functions.} The mean recursion of mode $k$ started at $(w,w')=(1,0)$ is $g_k$, so $e_1^\top B_k^\tau e_1=g_k(\tau)^2$ for $B_k=\widetilde B_k-\eta^2c_ke_1e_1^\top$, and $e_1^\top\widetilde B_k^\tau e_1=h_k(\tau)/\eta^2$ by expanding the rank-one perturbation. From $g_k(\tau)=\frac{r_+^{\tau+1}-r_-^{\tau+1}}{r_+-r_-}$ with $r_++r_-=a_k$, $r_+r_-=\beta$, summing the geometric series gives $\widehat\psi_k$; the three poles $1/\beta$, $r_\pm^{-2}$ are the reciprocals of the eigenvalues $\beta,r_\pm^2$ of $B_k$. At $z=1$, $\widehat\psi_k(1)=\frac{\eta}{2(1-\beta)\lambda_k(1-x_k)}$ and $\widehat h_k(1)=\frac{\eta}{2(1-\beta)\lambda_k(1-\alpha_kx_k)}$, because $(1-x_k)(1-c_k\widehat\psi_k(1))=1-\alpha_kx_k$. Hence $p\gamma\sum_k\lambda_k^2\widehat h_k(1)=S(\eta)$ whenever no pole lies at $1$.

\emph{(b).} From $Q_0=0$, $Q_t=\sum_{s<t}\T_0^{t-1-s}(\Delta(Q_s)+G)$ and Lemma~\ref{lem:V} give $x_{t,k}=\E w_{t,k}^2=\sum_{s<t}\psi_k(t-1-s)\big(\sigma_k^2+p\gamma\lambda_k\ell_s+c_kx_{s,k}\big)$, whose solution is $x_{t,k}=\sum_{s<t}h_k(t-1-s)(\sigma_k^2+p\gamma\lambda_k\ell_s)$; multiplying by $\lambda_k$ and summing gives the equation, and generating functions give $\widehat\ell=\frac{z}{1-z}\widehat\phi+z\widehat K\widehat\ell$.

\emph{(a).} The matrix $\mathcal M=\operatorname{diag}(\alpha_k\lambda_k)+p\kappa_\beta\gamma\lambda\mathbf 1^\top$ is similar to $\widetilde{\mathcal M}=\operatorname{diag}(\alpha_k\lambda_k)+p\kappa_\beta\gamma\sqrt\lambda\sqrt\lambda^\top$, whose diagonal entries $\lambda_k(1+p\kappa_\beta(\gamma(3+\kappa_k)-1))\ge\lambda_k$ and off-diagonal entries are positive. The determinant lemma, $\det(\mu-\widetilde{\mathcal M})=\prod_k(\mu-\alpha_k\lambda_k)\big(1-p\kappa_\beta\gamma\sum_k\frac{\lambda_k}{\mu-\alpha_k\lambda_k}\big)$, the Perron vector of $\widetilde{\mathcal M}+cI$ ($c$ large), which is positive and so excludes $\mu^\star=\alpha_k\lambda_k$, and the strict decrease of the sum in $\mu$ show that $\mu^\star=\lambda_{\max}(\widetilde{\mathcal M})$ is the unique root with $\mu>\max(0,\max_k\alpha_k\lambda_k)$; with $\mu=2(1+\beta)/\eta$ the sum equals $S(\eta)$. By Lemma~\ref{lem:V}, $\rho(\T)=\rho(\T_{\mathrm{diag}})$, and the stationary equations of $\T_{\mathrm{diag}}$ with $\Sigma_\xi$ replaced by $\operatorname{diag}(\sigma_k^2)$ ($\rho(\T)$ does not depend on $\Sigma_\xi$) are the diagonal entries of Theorem~\ref{thm:lyap}(2), i.e.\ $(I-N)c'=b$ for $c'_k=2\lambda_k\E w_k^2/(\eta(1+\beta))$, $N=\frac{\eta}{2(1+\beta)}\mathcal M\ge0$ and $b=\sigma^2/(1-\beta^2)>0$. If $\rho(\T)<1$, a stationary solution exists with $c'\ge b>0$ and $Nc'<c'$, so $\rho(N)<1$ by the Collatz--Wielandt formula, i.e.\ $\eta<\eta^\star$. Conversely, for $\eta<\eta^\star$, $I-N$ is a nonsingular M-matrix and $c'>0$; with $\Sigma_0=\operatorname{diag}(\E w_k^2)$ and $\Sigma_1=A\Sigma_0/(1+\beta)$ the matrix $Q=\begin{psmallmatrix}\Sigma_0&\Sigma_1\\\Sigma_1&\Sigma_0\end{psmallmatrix}$ is positive definite ($|1-2x_k|<1$ because $\eta<\eta^\star\le2(1+\beta)/\lambda_1$, as $\mu^\star\ge e_1^\top\widetilde{\mathcal M}e_1\ge\lambda_1$) and solves $Q=\T(Q)+G$; then $Q_\varepsilon=Q-\varepsilon G\succ0$ and $Q_\varepsilon-\T(Q_\varepsilon)=(1-\varepsilon)G+\varepsilon\T(G)\succ0$ for small $\varepsilon>0$, which forces $\rho(\T)<1$. The formula for $\ell_\infty$ follows from $\E w_k^2=\widehat h_k(1)(\sigma_k^2+p\gamma\lambda_k\ell_\infty)$.

\emph{(c).} Write $\T_{\mathrm{diag}}=\mathcal A+\eta^2p\gamma rr^\top$, $\mathcal A=\operatorname{blkdiag}(\widetilde B_k)$; then $\eta^2p\gamma\,r^\top(I-z\mathcal A)^{-1}r=\widehat K(z)$ and $\det(I-z\T_{\mathrm{diag}})=\prod_k\det(I-z\widetilde B_k)\,\big(1-z\widehat K(z)\big)$. $\T_{\mathrm{diag}}$ preserves the cone of matrices whose $P$ blocks $\begin{psmallmatrix}x_k&y_k\\y_k&m_k\end{psmallmatrix}$ are positive semidefinite, so $\rho=\rho(\T_{\mathrm{diag}})>0$ has an eigenvector $X\ne0$ in this cone. If $r^\top X=\sum_k\lambda_kx_k=0$, then every $x_k=0$, hence $y_k=0$, and the third coordinate of $\T_{\mathrm{diag}}X$ is $x_k=0=\rho m_k$, so $X=0$; thus $r^\top X>0$ and $(\rho-\mathcal A)X=\eta^2p\gamma(r^\top X)r\ne0$. If $\rho$ were an eigenvalue of $\widetilde B_k$, the $(1,1)$ cofactor of $\rho-\widetilde B_k$, which equals $\rho(\rho+\beta)>0$ for every $c_k$, would make the adjugate nonzero with a left null vector $\ell$, $\ell_1\ne0$, so $e_1\notin\operatorname{range}(\rho-\widetilde B_k)$, a contradiction. Hence $X=\eta^2p\gamma(r^\top X)(\rho-\mathcal A)^{-1}r$ and $z\widehat K(z)=1$ at $z=1/\rho$. Any positive root $z'$ of $z\widehat K=1$ at which $\widehat K$ is finite makes the determinant vanish, so $1/z'\le\rho$.

\emph{(d).} With $c_k\ge0$, $h_k\ge0$, $K_\tau\ge0$ and $K_0=p\gamma\eta^2\sum_k\lambda_k^2>0$. Each $\widehat h_k$ is rational with nonnegative coefficients, so its smallest-modulus singularity is a positive pole where it tends to $+\infty$ (Pringsheim), or it is a polynomial; hence $z\widehat K$ increases from $0$ to $+\infty$ on $[0,R_K)$, $R_K=\min_k R_k$, and has a unique root there, which is $z_\star=1/\rho$ by (c), and $z_\star>1$ since $\rho<1$. For $|z|\le z_\star$, $z\ne z_\star$: $\operatorname{Re}(z\widehat K(z))=\sum_\tau K_\tau\operatorname{Re}z^{\tau+1}<\sum_\tau K_\tau z_\star^{\tau+1}=1$, since equality would need $z^{\tau+1}>0$ for $\tau=0$. So $1-z\widehat K$ has no other zero in $|z|\le z_\star$; the zero is simple as $\frac{d}{dz}(z\widehat K)(z_\star)=\sum_\tau(\tau+1)K_\tau z_\star^\tau>0$; and the poles of $\widehat\phi$ lie in $|z|\ge R_K>z_\star$. The rational function $\widehat\ell$ thus has, in $|z|\le z_\star$, only simple poles at $1$ (residue $-\ell_\infty$) and at $z_\star$, which gives the expansion with the stated $C$; $C>0$ because $\widehat\phi(z_\star)>0$, $z_\star>1$ and the derivative is positive.

\emph{(e).} At $\eta^\star$, $\rho(\T)=1$ by continuity, so $z_\star=1$; and $1<R_K$: $c_k\ge0$ gives $\alpha_k\ge1$, the admissible range gives $\alpha_kx_k<1$ and hence $x_k<1$, so the poles of $\widehat\psi_k$ lie outside the closed unit disk and $c_kz\widehat\psi_k(z)\le c_k\widehat\psi_k(1)<1$ on $[0,1]$, the last inequality being $(1-x_k)(1-c_k\widehat\psi_k(1))=1-\alpha_kx_k>0$; hence $R_k>1$. The argument of (d) applies at $\eta^\star$ as well, so $z_\star(\eta)$ is the branch through $(1,\eta^\star)$. Thus $\Phi(z,\eta)=z\widehat K(z;\eta)$ is analytic near $(1,\eta^\star)$, with $\partial_z\Phi=\sum_\tau(\tau+1)K_\tau>0$ and $\partial_\eta\Phi=S'(\eta^\star)$ because $\widehat K(1;\eta)=S(\eta)$ there. Each term of $S$ has the form $\eta/(1-\eta d_k)$ with $\eta d_k<1$, so $S'>0$. The implicit function theorem gives $z_\star(\eta)$ with $z_\star'(\eta^\star)=-S'/\partial_z\Phi$, and $\rho=1/z_\star$.

\paragraph{Corollary~\ref{cor:special}.}
(i) With $\Lambda=\eta I$, $A=H$, $S=B$ and $\mu=\beta$, $\frac1B\tr[\Lambda AG_\mu^{-1}]=\frac{\eta}{B}\sum_k\frac{\lambda_k}{2(1-\beta)-(\frac{1-\beta}{1+\beta}+\frac1B)\eta\lambda_k}$, which is $S(\eta)$ at $\gamma=1$, $\kappa_k=0$ ($\alpha_k=1+p\kappa_\beta$); equation~(56) of \citet{wang2026momentum} is the same sum. (ii) Substitute $\lambda_k=q\tilde\lambda_k$, $\gamma=1/q$, $\kappa_k=0$, $\eta=(1-\beta)\eta_{\mathrm{EMA}}$ in $S$; for $\tilde\lambda_k\equiv1$ the equation is linear in $\eta_{\mathrm{EMA}}$, and in (c) all modes are equal, so $1=z\widehat K(z)$ becomes $Q_3(z)=\eta^2(Pp\gamma\lambda^2+c)z(1+\beta z)$ with $Q_3$ the cubic denominator of $\widehat\psi$. (iii) $\alpha_k\lambda_k=q_k\tilde\lambda_k+p\kappa_\beta\tilde\lambda_k(3+\tilde\kappa_k-2q_k)$.
(iv) From $\alpha_k\lambda_k=\lambda_k+\kappa_\beta c_k/\lambda_k$, $\alpha_kx_k=\eta_{\mathrm{EMA}}d_k$ and $\frac{\eta p\gamma}{2(1-\beta)}=\frac{\eta_{\mathrm{EMA}}p\gamma}2$; with $\mu^\star=\lambda_{\max}(\widetilde{\mathcal M})$ from the proof of Theorem~\ref{thm:main}(a), $1/\eta^\star_{\mathrm{EMA}}=\mu^\star/(2\kappa_\beta)$ gives the second formula, and $\widetilde{\mathcal M}=H+p\kappa_\beta M_\infty$ gives (c). (a) With $y=1/\eta_{\mathrm{EMA}}$, the level $s$ is the root of $F_\theta(y)=\sum_k\lambda_k/(y-d_k)=2s/(p\gamma)$ on $y>\max(\max_kd_k,0)$, where $F_\theta$ decreases; it exists for $s\le1$ because it exists for $s=1$. From $\beta=0$ to $\beta$ every $d_k$ decreases by $\delta_k=\frac{\beta}{1+\beta}\lambda_k$. If $y_0$ is the root at $\beta=0$, then $y_0\ge\mu^\star/2\ge\lambda_{\max}/2>\max_k\delta_k$, $F_\theta(y_0-\min_k\delta_k)\le F_1(y_0)$ and $F_\theta(y_0-\max_k\delta_k)\ge F_1(y_0)$ (when $y_0-\max_k\delta_k$ is admissible), so the root at $\beta$ lies in $[y_0-\max_k\delta_k,\,y_0-\min_k\delta_k]$; the same comparison between $\beta_1<\beta_2$ gives the strict increase. (b) As $\theta\downarrow0$, $\operatorname{diag}(d_k)$ decreases to $\operatorname{diag}(c_k/(2\lambda_k))=\frac p2\operatorname{diag}((\gamma(2+\kappa_k)-1)\lambda_k)$. Since $(\gamma(2+\kappa_k)-1)\lambda_k\ge-\lambda_k$ and $\gamma\ge1$, $M_\infty\succeq-H+\sqrt\lambda\sqrt\lambda^\top$, whose principal $2\times2$ submatrix on the two largest eigenvalues has top eigenvalue $\sqrt{\lambda_{(1)}\lambda_{(2)}}$; Cauchy interlacing gives the bound. For $s<1$ the limiting root is finite: a mode with $c_k>0$ gives a pole, a mode with $c_k=0$ an unbounded term, and if every $c_k<0$ the limiting sum tends to $\sum_k\gamma/(1-\gamma(2+\kappa_k))\ge P\gamma\ge1>s$. (c) Since $G\succ0$, $\lambda_{\max}(H+\tau M_\infty)<2\lambda_{\max}$ if and only if $\tau G^{-1/2}M_\infty G^{-1/2}\prec I$; with $\tau=\frac{1+\beta}{(1-\beta)B}$ this is $B>B_{1/2}(\beta)$, and equality on one side is equality on the other. Ostrowski's theorem with $\lambda_{\max}I\preceq G\preceq2\lambda_{\max}I$ gives the bracket.

\paragraph{Proposition~\ref{prop:1d}.}
Here $\widetilde{\mathcal M}=\alpha H+p\kappa_\beta\sqrt\lambda\sqrt\lambda^\top$ in the eigenbasis of $H$ and $\mu^\star=\lambda_{\max}(\widetilde{\mathcal M})$. The Rayleigh quotients at $e_1$ and at $\sqrt\lambda/\sqrt{\tr H}$ give $\mu^\star\ge(\alpha+p\kappa_\beta)\lambda_1=\mu_{\mathrm{1D}}$ and $\mu^\star\ge\alpha\tr H^2/\tr H+p\kappa_\beta\tr H$; Weyl's inequality gives $\mu^\star\le\max_k\alpha\lambda_k+p\kappa_\beta\tr H\le\max(\alpha,0)\lambda_1+p\kappa_\beta\tr H$. The ratio of thresholds is $\mu^\star/\mu_{\mathrm{1D}}$.

\paragraph{Proposition~\ref{prop:floor}.}
Here $\sigma_k^2=\sigma^2p\lambda_k$, so $f_\infty=\frac{\sigma^2}{\gamma}S$ and Theorem~\ref{thm:main}(a) gives $\ell_\infty$. With $c_k\ge0$, $K_\tau\ge0$, $\sum_\tau K_\tau=\widehat K(1)=S$ and $\sum_\tau K_\tau z_\star^{\tau+1}=1$ with $z_\star=1/\rho(\T)>1$; by convexity of $x\mapsto z_\star^x$, $1/S=\sum_\tau\frac{K_\tau}{S}z_\star^{\tau+1}\ge z_\star^{\bar\tau}$, so $\rho(\T)\ge S^{1/\bar\tau}$, and $\bar\tau\ge1$, $S<1$ give $S^{1/\bar\tau}\ge S$. The table is part~G of \texttt{note-v2-checks.py}.

\paragraph{Proposition~\ref{prop:eff}.}
With $\eta=(1-\beta)\eta_{\mathrm{eff}}$, Theorem~\ref{thm:lyap}(2) reads $M_\beta\Sigma_0=\eta_{\mathrm{eff}}\Sigma_\xi$, $M_\beta=\mathcal L_0-\eta_{\mathrm{eff}}\mathcal A_\beta$, $\mathcal A_\beta=(1-\beta)\mathcal B_\beta$, and $M_\beta\succ0$ at stable $\beta$ by Proposition~\ref{prop:cai}. In the eigenbasis of $H$, $\mathcal A_0-\mathcal A_\beta$ multiplies entry $(j,k)$ by $m_{jk}(\beta)=\frac{2\beta}{1+\beta}\lambda_j\lambda_k+\frac{\beta}{1-\beta^2}(\lambda_j-\lambda_k)^2$, which increases strictly in $\beta$ and is positive. (i) By Proposition~\ref{prop:cai}, $1/\eta^\star_{\mathrm{eff}}(\beta)=\lambda_{\max}(\mathcal L_0^{-1/2}\mathcal A_\beta\mathcal L_0^{-1/2})$, and $\mathcal L_0^{-1/2}(\mathcal A_{\beta_1}-\mathcal A_{\beta_2})\mathcal L_0^{-1/2}\succ0$ for $\beta_1<\beta_2$, so Weyl's inequality gives the strict increase; $m_{jk}(\beta)\ge\frac{\beta}{1+\beta}(\lambda_j^2+\lambda_k^2)\ge\frac{\beta}{1+\beta}\lambda_{\min}(H)(\lambda_j+\lambda_k)$ gives $\mathcal L_0^{-1/2}(\mathcal A_0-\mathcal A_\beta)\mathcal L_0^{-1/2}\succeq\frac{\beta}{1+\beta}\lambda_{\min}(H)I$ and the margin. (ii) $M_{\beta_2}-M_{\beta_1}\succ0$, $M_{\beta_2}^{-1}\prec M_{\beta_1}^{-1}$, and $\tr(\Sigma_\xi\Sigma_0)=\eta_{\mathrm{eff}}\langle\Sigma_\xi,M_\beta^{-1}\Sigma_\xi\rangle$ decreases strictly. The example is computed in \texttt{note-v2-checks.py} (part E), where both values of $\beta$ have $\rho(\T)<1$ (0.809 and 0.966); the minibatch Hessian $H+N$ is indefinite (eigenvalues $-0.466$, $0.966$) and $H-N$ is positive definite.

\renewcommand{\bibfont}{\footnotesize}\setlength{\bibsep}{0pt}
\bibliographystyle{plainnat}
\bibliography{refs-v2}

\end{document}
